% Requires in the preamble:
% \usepackage{amsmath,amssymb}

\chapter{Tables}
\label{app:tables}

A table is a map.  It is not the trip.

The formulas in this appendix are meant for quick reference after the ideas have been
learned.  Before using a formula, ask three questions.

\[
\text{What does the formula mean?}
\]

\[
\text{What hypotheses does it require?}
\]

\[
\text{Does the answer make sense in this problem?}
\]

A derivative formula should agree with the idea of local rate of change.  An integral
formula should agree with accumulated signed amount.  A series test should be used
only when its hypotheses are satisfied.  A numerical rule should come with an error
estimate when accuracy matters.

Unless degrees are explicitly mentioned, all trigonometric formulas in calculus use
radians.

\section{Derivative formulas}
\label{appsec:derivative-formulas}

The derivative measures local change.  These rules tell how local change behaves when
functions are built from simpler functions.

\subsection{General differentiation rules}
\label{appsubsec:general-differentiation-rules}

In the table below, \(c\) is a constant and all functions are assumed differentiable
where the formula is used.

\begin{center}
\renewcommand{\arraystretch}{1.8}
\begin{tabular}{p{0.32\textwidth}|p{0.52\textwidth}}
\textbf{rule} & \textbf{formula} \\ \hline

Constant rule &
\(\displaystyle \frac{d}{dx}(c)=0\) \\

Constant multiple rule &
\(\displaystyle \frac{d}{dx}\bigl(c f(x)\bigr)=c f'(x)\) \\

Sum rule &
\(\displaystyle \frac{d}{dx}\bigl(f(x)+g(x)\bigr)=f'(x)+g'(x)\) \\

Difference rule &
\(\displaystyle \frac{d}{dx}\bigl(f(x)-g(x)\bigr)=f'(x)-g'(x)\) \\

Product rule &
\(\displaystyle \frac{d}{dx}\bigl(f(x)g(x)\bigr)=f'(x)g(x)+f(x)g'(x)\) \\

Quotient rule &
\(\displaystyle
\frac{d}{dx}\left(\frac{f(x)}{g(x)}\right)
=
\frac{g(x)f'(x)-f(x)g'(x)}{(g(x))^2}
\),
where \(g(x)\neq0\) \\

Chain rule &
\(\displaystyle
\frac{d}{dx}f(g(x))=f'(g(x))g'(x)
\) \\

Inverse function rule &
If \(y=f^{-1}(x)\), then
\(\displaystyle
\frac{d}{dx}f^{-1}(x)=\frac1{f'(y)}
\),
provided \(f'(y)\neq0\) \\

Implicit differentiation &
Differentiate both sides with respect to \(x\), remembering that \(y=y(x)\), so
\(\displaystyle \frac{d}{dx}y=y'\) \\

Logarithmic differentiation &
For positive \(y=f(x)\), take logarithms first:
\(\displaystyle \ln y=\ln f(x)\), then differentiate implicitly.
\end{tabular}
\end{center}

\begin{quote}
\textbf{Warning.}
The derivative is linear, but it is not multiplicative:

\[
(fg)'\neq f'g'
\]

in general.  The missing terms are exactly what the product rule supplies.
\end{quote}

\subsection{Elementary derivatives}
\label{appsubsec:elementary-derivatives}

\begin{center}
\renewcommand{\arraystretch}{1.9}
\begin{tabular}{p{0.40\textwidth}|p{0.44\textwidth}}
\textbf{function} & \textbf{derivative} \\ \hline

\(\displaystyle x^r\) &
\(\displaystyle r x^{r-1}\), on intervals where \(x^r\) is defined and differentiable \\

\(\displaystyle e^x\) &
\(\displaystyle e^x\) \\

\(\displaystyle e^{kx}\) &
\(\displaystyle k e^{kx}\) \\

\(\displaystyle a^x\), \(a>0\) &
\(\displaystyle a^x\ln a\) \\

\(\displaystyle \ln x\), \(x>0\) &
\(\displaystyle \frac1x\) \\

\(\displaystyle \ln|x|\), \(x\neq0\) &
\(\displaystyle \frac1x\) \\

\(\displaystyle \log_a x\), \(a>0,\ a\neq1,\ x>0\) &
\(\displaystyle \frac1{x\ln a}\)
\end{tabular}
\end{center}

\subsection{Trigonometric derivatives}
\label{appsubsec:trigonometric-derivatives}

\begin{center}
\renewcommand{\arraystretch}{1.8}
\begin{tabular}{c|c}
\textbf{function} & \textbf{derivative} \\ \hline
\(\sin x\) & \(\cos x\) \\
\(\cos x\) & \(-\sin x\) \\
\(\tan x\) & \(\sec^2 x\) \\
\(\cot x\) & \(-\csc^2 x\) \\
\(\sec x\) & \(\sec x\tan x\) \\
\(\csc x\) & \(-\csc x\cot x\)
\end{tabular}
\end{center}

These formulas require radian measure.  For example,

\[
\frac{d}{dx}\sin x=\cos x
\]

is a radian formula.

\subsection{Inverse trigonometric derivatives}
\label{appsubsec:inverse-trig-derivatives}

\begin{center}
\renewcommand{\arraystretch}{2.0}
\begin{tabular}{p{0.34\textwidth}|p{0.34\textwidth}|p{0.18\textwidth}}
\textbf{function} & \textbf{derivative} & \textbf{where} \\ \hline

\(\arcsin x\) &
\(\displaystyle \frac1{\sqrt{1-x^2}}\) &
\(-1<x<1\) \\

\(\arccos x\) &
\(\displaystyle -\frac1{\sqrt{1-x^2}}\) &
\(-1<x<1\) \\

\(\arctan x\) &
\(\displaystyle \frac1{1+x^2}\) &
all real \(x\) \\

\(\operatorname{arcsec} x\) &
\(\displaystyle \frac1{|x|\sqrt{x^2-1}}\) &
\(|x|>1\) \\

\(\operatorname{arccsc} x\) &
\(\displaystyle -\frac1{|x|\sqrt{x^2-1}}\) &
\(|x|>1\)
\end{tabular}
\end{center}

Different texts use slightly different conventions for \(\operatorname{arccot}\),
\(\operatorname{arcsec}\), and \(\operatorname{arccsc}\).  When using a calculator or
computer algebra system, check its convention.

\subsection{Second derivative reminders}
\label{appsubsec:second-derivative-reminders}

If \(f\) is twice differentiable, then

\[
f''(x)=\frac{d}{dx}f'(x).
\]

The signs of \(f'\) and \(f''\) have different meanings.

\begin{center}
\renewcommand{\arraystretch}{1.6}
\begin{tabular}{c|c}
\textbf{condition} & \textbf{graph meaning} \\ \hline
\(f'(x)>0\) & \(f\) is increasing there \\
\(f'(x)<0\) & \(f\) is decreasing there \\
\(f''(x)>0\) & graph is concave up there \\
\(f''(x)<0\) & graph is concave down there
\end{tabular}
\end{center}

A zero derivative does not automatically mean a maximum or minimum.  A zero second
derivative does not automatically mean an inflection point.  Use sign changes or other
information.

\section{Integral formulas}
\label{appsec:integral-formulas}

An indefinite integral is an antiderivative family:

\[
\int f(x)\,dx=F(x)+C
\]

means

\[
F'(x)=f(x).
\]

A definite integral is accumulated signed amount:

\[
\int_a^b f(x)\,dx.
\]

The two are connected by the Fundamental Theorem of Calculus.

\subsection{General integral rules}
\label{appsubsec:general-integral-rules}

\begin{center}
\renewcommand{\arraystretch}{1.9}
\begin{tabular}{p{0.32\textwidth}|p{0.52\textwidth}}
\textbf{rule} & \textbf{formula} \\ \hline

Linearity &
\(\displaystyle
\int \bigl(af(x)+bg(x)\bigr)\,dx
=
a\int f(x)\,dx+b\int g(x)\,dx
\) \\

Definite linearity &
\(\displaystyle
\int_a^b \bigl(af(x)+bg(x)\bigr)\,dx
=
a\int_a^b f(x)\,dx+b\int_a^b g(x)\,dx
\) \\

Orientation &
\(\displaystyle
\int_b^a f(x)\,dx=-\int_a^b f(x)\,dx
\) \\

Zero interval &
\(\displaystyle
\int_a^a f(x)\,dx=0
\) \\

Additivity over intervals &
\(\displaystyle
\int_a^b f(x)\,dx+\int_b^c f(x)\,dx=\int_a^c f(x)\,dx
\) \\

Fundamental Theorem &
If \(F'=f\), then
\(\displaystyle
\int_a^b f(x)\,dx=F(b)-F(a)
\) \\

Substitution &
If \(u=g(x)\), then
\(\displaystyle
\int f(g(x))g'(x)\,dx=\int f(u)\,du
\) \\

Definite substitution &
If \(u=g(x)\), then
\(\displaystyle
\int_a^b f(g(x))g'(x)\,dx
=
\int_{g(a)}^{g(b)} f(u)\,du
\) \\

Integration by parts &
\(\displaystyle
\int u\,dv=uv-\int v\,du
\) \\

Definite integration by parts &
\(\displaystyle
\int_a^b u\,dv=\bigl[uv\bigr]_a^b-\int_a^b v\,du
\)
\end{tabular}
\end{center}

\begin{quote}
\textbf{Check by differentiating.}
For an indefinite integral, the fastest way to check an answer is to differentiate it.
The derivative should return the original integrand.
\end{quote}

\subsection{Basic antiderivatives}
\label{appsubsec:basic-antiderivatives}

\begin{center}
\renewcommand{\arraystretch}{2.0}
\begin{tabular}{p{0.42\textwidth}|p{0.42\textwidth}}
\textbf{integrand} & \textbf{antiderivative} \\ \hline

\(\displaystyle x^r,\quad r\neq -1\) &
\(\displaystyle \int x^r\,dx=\frac{x^{r+1}}{r+1}+C\) \\

\(\displaystyle \frac1x\) &
\(\displaystyle \int \frac1x\,dx=\ln|x|+C\) \\

\(\displaystyle e^x\) &
\(\displaystyle \int e^x\,dx=e^x+C\) \\

\(\displaystyle e^{kx}\), \(k\neq0\) &
\(\displaystyle \int e^{kx}\,dx=\frac1k e^{kx}+C\) \\

\(\displaystyle a^x,\quad a>0,\ a\neq1\) &
\(\displaystyle \int a^x\,dx=\frac{a^x}{\ln a}+C\) \\

\(\displaystyle \frac{f'(x)}{f(x)}\) &
\(\displaystyle \int \frac{f'(x)}{f(x)}\,dx=\ln|f(x)|+C\)
\end{tabular}
\end{center}

\subsection{Trigonometric antiderivatives}
\label{appsubsec:trig-antiderivatives}

\begin{center}
\renewcommand{\arraystretch}{1.8}
\begin{tabular}{c|c}
\textbf{integrand} & \textbf{antiderivative} \\ \hline
\(\sin x\) & \(-\cos x+C\) \\
\(\cos x\) & \(\sin x+C\) \\
\(\sec^2 x\) & \(\tan x+C\) \\
\(\csc^2 x\) & \(-\cot x+C\) \\
\(\sec x\tan x\) & \(\sec x+C\) \\
\(\csc x\cot x\) & \(-\csc x+C\) \\
\(\tan x\) & \(-\ln|\cos x|+C\) \\
\(\cot x\) & \(\ln|\sin x|+C\) \\
\(\sec x\) & \(\ln|\sec x+\tan x|+C\) \\
\(\csc x\) & \(\ln|\csc x-\cot x|+C\)
\end{tabular}
\end{center}

For constants \(k\neq0\),

\[
\int \sin(kx)\,dx=-\frac1k\cos(kx)+C,
\]

and

\[
\int \cos(kx)\,dx=\frac1k\sin(kx)+C.
\]

\subsection{Inverse trigonometric forms}
\label{appsubsec:inverse-trig-integral-forms}

For \(a>0\), the following patterns are especially common.

\begin{center}
\renewcommand{\arraystretch}{2.1}
\begin{tabular}{p{0.42\textwidth}|p{0.42\textwidth}}
\textbf{integrand} & \textbf{antiderivative} \\ \hline

\(\displaystyle \frac1{a^2+x^2}\) &
\(\displaystyle \frac1a\arctan\left(\frac{x}{a}\right)+C\) \\

\(\displaystyle \frac1{\sqrt{a^2-x^2}}\) &
\(\displaystyle \arcsin\left(\frac{x}{a}\right)+C\) \\

\(\displaystyle \frac1{x^2-a^2}\) &
\(\displaystyle
\frac1{2a}\ln\left|\frac{x-a}{x+a}\right|+C
\) \\

\(\displaystyle \frac1{a^2-x^2}\) &
\(\displaystyle
\frac1{2a}\ln\left|\frac{a+x}{a-x}\right|+C
\) \\

\(\displaystyle \frac1{\sqrt{x^2+a^2}}\) &
\(\displaystyle
\ln\left|x+\sqrt{x^2+a^2}\right|+C
\) \\

\(\displaystyle \frac1{\sqrt{x^2-a^2}}\) &
\(\displaystyle
\ln\left|x+\sqrt{x^2-a^2}\right|+C
\)
\end{tabular}
\end{center}

These formulas often appear after completing the square.

\subsection{Trigonometric substitution patterns}
\label{appsubsec:trig-substitution-patterns}

For \(a>0\), the following substitutions match the Pythagorean identities.

\begin{center}
\renewcommand{\arraystretch}{1.8}
\begin{tabular}{p{0.28\textwidth}|p{0.24\textwidth}|p{0.32\textwidth}}
\textbf{expression} & \textbf{substitution} & \textbf{identity used} \\ \hline

\(\sqrt{a^2-x^2}\) &
\(x=a\sin\theta\) &
\(1-\sin^2\theta=\cos^2\theta\) \\

\(\sqrt{a^2+x^2}\) &
\(x=a\tan\theta\) &
\(1+\tan^2\theta=\sec^2\theta\) \\

\(\sqrt{x^2-a^2}\) &
\(x=a\sec\theta\) &
\(\sec^2\theta-1=\tan^2\theta\)
\end{tabular}
\end{center}

After using a trigonometric substitution, convert the answer back to \(x\) unless the
problem asks for a trigonometric parameter.

\subsection{Partial fraction patterns}
\label{appsubsec:partial-fraction-patterns}

Partial fractions are used for rational functions.  Factor the denominator first.

\begin{center}
\renewcommand{\arraystretch}{1.9}
\begin{tabular}{p{0.32\textwidth}|p{0.52\textwidth}}
\textbf{denominator factor} & \textbf{partial fraction form} \\ \hline

Distinct linear factor \(x-a\) &
\(\displaystyle \frac{A}{x-a}\) \\

Repeated linear factor \((x-a)^m\) &
\(\displaystyle
\frac{A_1}{x-a}
+
\frac{A_2}{(x-a)^2}
+\cdots+
\frac{A_m}{(x-a)^m}
\) \\

Irreducible quadratic factor \(ax^2+bx+c\) &
\(\displaystyle
\frac{Ax+B}{ax^2+bx+c}
\) \\

Repeated irreducible quadratic factor \((ax^2+bx+c)^m\) &
Use one numerator \(A_jx+B_j\) over each power
\[
(ax^2+bx+c)^j,
\qquad j=1,\ldots,m.
\]
\end{tabular}
\end{center}

If the degree of the numerator is at least the degree of the denominator, divide first.

\section{Common Taylor series}
\label{appsec:common-taylor-series}

Taylor polynomials approximate functions near a chosen center.  Taylor series continue
that approximation indefinitely.

\subsection{Taylor polynomial and remainder}
\label{appsubsec:taylor-polynomial-remainder-table}

The Taylor polynomial of degree \(N\) for \(f\) centered at \(a\) is

\[
\boxed{
T_N(x)
=
\sum_{n=0}^{N}\frac{f^{(n)}(a)}{n!}(x-a)^n.
}
\]

That is,

\[
T_N(x)
=
f(a)+f'(a)(x-a)+\frac{f''(a)}{2!}(x-a)^2+\cdots+
\frac{f^{(N)}(a)}{N!}(x-a)^N.
\]

Taylor's theorem with Lagrange remainder says that, under the usual differentiability
hypotheses,

\[
\boxed{
f(x)=T_N(x)+R_N(x),
}
\]

where

\[
\boxed{
R_N(x)
=
\frac{f^{(N+1)}(\xi)}{(N+1)!}(x-a)^{N+1}
}
\]

for some \(\xi\) between \(a\) and \(x\).

Thus, if

\[
|f^{(N+1)}(t)|\le M
\]

between \(a\) and \(x\), then

\[
\boxed{
|R_N(x)|\le
\frac{M}{(N+1)!}|x-a|^{N+1}.
}
\]

This is the formula that turns a Taylor polynomial into a controlled approximation.

\subsection{Maclaurin series}
\label{appsubsec:maclaurin-series-table}

A Taylor series centered at \(0\) is called a Maclaurin series.

\begin{center}
\small
\renewcommand{\arraystretch}{2.0}
\begin{tabular}{p{0.28\textwidth}|p{0.42\textwidth}|p{0.18\textwidth}}
\textbf{function} & \textbf{series} & \textbf{converges for} \\ \hline

\(\displaystyle e^x\) &
\(\displaystyle
\sum_{n=0}^{\infty}\frac{x^n}{n!}
=
1+x+\frac{x^2}{2!}+\frac{x^3}{3!}+\cdots
\) &
all \(x\) \\

\(\displaystyle \sin x\) &
\(\displaystyle
\sum_{n=0}^{\infty}(-1)^n\frac{x^{2n+1}}{(2n+1)!}
=
x-\frac{x^3}{3!}+\frac{x^5}{5!}-\cdots
\) &
all \(x\) \\

\(\displaystyle \cos x\) &
\(\displaystyle
\sum_{n=0}^{\infty}(-1)^n\frac{x^{2n}}{(2n)!}
=
1-\frac{x^2}{2!}+\frac{x^4}{4!}-\cdots
\) &
all \(x\) \\

\(\displaystyle \frac1{1-x}\) &
\(\displaystyle
\sum_{n=0}^{\infty}x^n
=
1+x+x^2+x^3+\cdots
\) &
\(|x|<1\) \\

\(\displaystyle \frac1{1+x}\) &
\(\displaystyle
\sum_{n=0}^{\infty}(-1)^n x^n
=
1-x+x^2-x^3+\cdots
\) &
\(|x|<1\) \\

\(\displaystyle \ln(1+x)\) &
\(\displaystyle
\sum_{n=1}^{\infty}(-1)^{n+1}\frac{x^n}{n}
=
x-\frac{x^2}{2}+\frac{x^3}{3}-\cdots
\) &
\(-1<x\le1\) \\

\(\displaystyle \ln(1-x)\) &
\(\displaystyle
-\sum_{n=1}^{\infty}\frac{x^n}{n}
=
-x-\frac{x^2}{2}-\frac{x^3}{3}-\cdots
\) &
\(-1\le x<1\) \\

\(\displaystyle \arctan x\) &
\(\displaystyle
\sum_{n=0}^{\infty}(-1)^n\frac{x^{2n+1}}{2n+1}
=
x-\frac{x^3}{3}+\frac{x^5}{5}-\cdots
\) &
\(-1\le x\le1\) \\

\(\displaystyle (1+x)^p\) &
\(\displaystyle
1+px+\frac{p(p-1)}{2!}x^2+
\frac{p(p-1)(p-2)}{3!}x^3+\cdots
\) &
\(|x|<1\), endpoints depend on \(p\)
\end{tabular}
\end{center}

\begin{quote}
\textbf{Radius warning.}
Inside the radius of convergence, power series may be differentiated and integrated
term by term.  At endpoints, test separately.
\end{quote}

\subsection{Series shifted from zero}
\label{appsubsec:shifted-series-table}

A series centered at \(a\) is obtained by replacing \(x\) with \(x-a\).

For example,

\[
e^{x}
=
e^a e^{x-a}.
\]

So

\[
\boxed{
e^x
=
e^a
\sum_{n=0}^{\infty}\frac{(x-a)^n}{n!}.
}
\]

Also,

\[
\frac1{1-x}
=
\frac1{1-a-(x-a)}
=
\frac1{1-a}\cdot
\frac1{1-\frac{x-a}{1-a}},
\]

so, for

\[
\left|\frac{x-a}{1-a}\right|<1,
\]

we have

\[
\boxed{
\frac1{1-x}
=
\sum_{n=0}^{\infty}\frac{(x-a)^n}{(1-a)^{n+1}}.
}
\]

Shifted series are useful when the function must be approximated near a point other
than \(0\).

\section{Series tests}
\label{appsec:series-tests}

A series

\[
\sum_{n=1}^{\infty}a_n
\]

converges if its partial sums approach a finite number.

The tests below are tools for deciding convergence.  They do not all give the sum.

\subsection{Basic comparison and structure tests}
\label{appsubsec:basic-series-tests-table}

\begin{center}
\small
\renewcommand{\arraystretch}{2.0}
\begin{tabular}{p{0.20\textwidth}|p{0.32\textwidth}|p{0.34\textwidth}}
\textbf{test} & \textbf{hypotheses/check} & \textbf{conclusion} \\ \hline

Term test for divergence &
Compute \(\displaystyle \lim_{n\to\infty}a_n\). &
If the limit is not \(0\), or does not exist, then \(\sum a_n\) diverges.
If the limit is \(0\), the test is inconclusive. \\

Geometric series &
\(\displaystyle \sum_{n=0}^{\infty} ar^n\). &
Converges if \(|r|<1\), with sum
\(\displaystyle \frac{a}{1-r}\).
Diverges if \(|r|\ge1\). \\

\(p\)-series &
\(\displaystyle \sum_{n=1}^{\infty}\frac1{n^p}\). &
Converges if \(p>1\).  Diverges if \(p\le1\). \\

Telescoping series &
Partial sums collapse after cancellation. &
Write \(S_N\), cancel, then take \(\displaystyle \lim_{N\to\infty}S_N\). \\

Direct comparison &
Assume \(0\le a_n\le b_n\) eventually. &
If \(\sum b_n\) converges, then \(\sum a_n\) converges.
If \(\sum a_n\) diverges, then \(\sum b_n\) diverges. \\

Limit comparison &
Assume \(a_n,b_n>0\) eventually and compute
\(\displaystyle L=\lim_{n\to\infty}\frac{a_n}{b_n}\). &
If \(0<L<\infty\), then \(\sum a_n\) and \(\sum b_n\) either both converge
or both diverge.
\end{tabular}
\end{center}

\subsection{Integral, alternating, ratio, and root tests}
\label{appsubsec:advanced-series-tests-table}

\begin{center}
\small
\renewcommand{\arraystretch}{2.0}
\begin{tabular}{p{0.20\textwidth}|p{0.34\textwidth}|p{0.32\textwidth}}
\textbf{test} & \textbf{hypotheses/check} & \textbf{conclusion} \\ \hline

Integral test &
\(a_n=f(n)\), where \(f\) is positive, continuous, and decreasing for large \(x\). &
\(\displaystyle \sum a_n\) and \(\displaystyle \int^\infty f(x)\,dx\) either both
converge or both diverge. \\

Integral-test remainder &
Same hypotheses as the integral test. &
For \(S=\sum_{n=1}^{\infty}a_n\) and \(S_N=\sum_{n=1}^{N}a_n\),

\[
\int_{N+1}^{\infty}f(x)\,dx
\le
S-S_N
\le
\int_N^{\infty}f(x)\,dx.
\] \\

Alternating series test &
\(\displaystyle \sum (-1)^{n-1}b_n\), where \(b_n\ge0\), \(b_n\) decreases, and
\(b_n\to0\). &
The series converges.  Also,

\[
|S-S_N|\le b_{N+1}.
\] \\

Absolute convergence &
Check whether \(\displaystyle \sum |a_n|\) converges. &
If \(\sum |a_n|\) converges, then \(\sum a_n\) converges. \\

Conditional convergence &
\(\displaystyle \sum a_n\) converges but \(\displaystyle \sum |a_n|\) diverges. &
The series converges conditionally.  Alternating series often behave this way. \\

Ratio test &
Compute
\[
L=\lim_{n\to\infty}\left|\frac{a_{n+1}}{a_n}\right|.
\] &
If \(L<1\), the series converges absolutely.
If \(L>1\) or \(L=\infty\), it diverges.
If \(L=1\), inconclusive. \\

Root test &
Compute
\[
L=\lim_{n\to\infty}\sqrt[n]{|a_n|}.
\] &
If \(L<1\), the series converges absolutely.
If \(L>1\) or \(L=\infty\), it diverges.
If \(L=1\), inconclusive.
\end{tabular}
\end{center}

\subsection{Power series convergence}
\label{appsubsec:power-series-tests-table}

For a power series

\[
\sum_{n=0}^{\infty} c_n(x-a)^n,
\]

the convergence set is centered at \(a\).

\begin{center}
\renewcommand{\arraystretch}{1.7}
\begin{tabular}{p{0.30\textwidth}|p{0.54\textwidth}}
\textbf{quantity} & \textbf{meaning} \\ \hline

Center &
The number \(a\) in powers of \((x-a)\). \\

Radius of convergence \(R\) &
The series converges absolutely for \(|x-a|<R\) and diverges for \(|x-a|>R\). \\

Endpoints &
The points \(a-R\) and \(a+R\).  Test them separately. \\

Interval of convergence &
The full set of \(x\)-values where the series converges, including any endpoints that
pass their tests.
\end{tabular}
\end{center}

The ratio test often gives the radius.  If

\[
L=\lim_{n\to\infty}
\left|
\frac{c_{n+1}(x-a)^{n+1}}{c_n(x-a)^n}
\right|
\]

has the form

\[
L=\frac{|x-a|}{R},
\]

then the radius of convergence is \(R\).

\begin{quote}
\textbf{Endpoint warning.}
The ratio and root tests usually decide what happens inside and outside the radius.
They usually do not decide what happens at the endpoints.
\end{quote}

\section{Numerical integration rules}
\label{appsec:numerical-integration-rules}

Numerical integration approximates

\[
I=\int_a^b f(x)\,dx
\]

by replacing the graph with simple shapes.

Let

\[
\Delta x=\frac{b-a}{n},
\qquad
x_j=a+j\Delta x,
\qquad
j=0,1,\ldots,n.
\]

The midpoint of the \(j\)-th subinterval \([x_{j-1},x_j]\) is

\[
m_j=\frac{x_{j-1}+x_j}{2}.
\]

\subsection{Rectangle, midpoint, trapezoidal, and Simpson rules}
\label{appsubsec:numerical-rules-table}

\begin{center}
\small
\renewcommand{\arraystretch}{2.1}
\begin{tabular}{p{0.18\textwidth}|p{0.52\textwidth}|p{0.18\textwidth}}
\textbf{rule} & \textbf{formula} & \textbf{geometry} \\ \hline

Left rectangle rule &
\(\displaystyle
L_n=\Delta x\sum_{j=0}^{n-1} f(x_j)
\) &
Left endpoint rectangles \\

Right rectangle rule &
\(\displaystyle
R_n=\Delta x\sum_{j=1}^{n} f(x_j)
\) &
Right endpoint rectangles \\

Midpoint rule &
\(\displaystyle
M_n=\Delta x\sum_{j=1}^{n} f(m_j)
\) &
Midpoint rectangles \\

Trapezoidal rule &
\(\displaystyle
T_n=
\frac{\Delta x}{2}
\left[
f(x_0)+2\sum_{j=1}^{n-1}f(x_j)+f(x_n)
\right]
\) &
Trapezoids \\

Simpson's rule &
For even \(n\),

\[
S_n=
\frac{\Delta x}{3}
\left[
f(x_0)
+
4\sum_{\substack{1\le j\le n-1\\ j\text{ odd}}} f(x_j)
+
2\sum_{\substack{2\le j\le n-2\\ j\text{ even}}} f(x_j)
+
f(x_n)
\right].
\] &
Parabolic arcs
\end{tabular}
\end{center}

\begin{quote}
\textbf{Simpson's rule requires even \(n\).}
The subintervals are paired, because each pair is approximated by a parabola.
\end{quote}

\subsection{Error bounds}
\label{appsubsec:numerical-error-bounds-table}

The following bounds assume \(f\) has the required derivatives on \([a,b]\).

\begin{center}
\small
\renewcommand{\arraystretch}{2.0}
\begin{tabular}{p{0.20\textwidth}|p{0.38\textwidth}|p{0.28\textwidth}}
\textbf{rule} & \textbf{hypothesis} & \textbf{error bound} \\ \hline

Left or right rectangle &
\(\displaystyle |f'(x)|\le K\) on \([a,b]\) &
\(\displaystyle
|I-L_n|\le \frac{K(b-a)^2}{2n},
\qquad
|I-R_n|\le \frac{K(b-a)^2}{2n}
\) \\

Midpoint &
\(\displaystyle |f''(x)|\le K\) on \([a,b]\) &
\(\displaystyle
|I-M_n|\le \frac{K(b-a)^3}{24n^2}
\) \\

Trapezoidal &
\(\displaystyle |f''(x)|\le K\) on \([a,b]\) &
\(\displaystyle
|I-T_n|\le \frac{K(b-a)^3}{12n^2}
\) \\

Simpson &
\(\displaystyle |f^{(4)}(x)|\le K\) on \([a,b]\), and \(n\) is even &
\(\displaystyle
|I-S_n|\le \frac{K(b-a)^5}{180n^4}
\)
\end{tabular}
\end{center}

These are worst-case guarantees.  The actual error may be much smaller.

\subsection{Monotonicity and concavity checks}
\label{appsubsec:monotonicity-concavity-numerical-checks}

Before using a numerical value, a quick shape check can tell whether the approximation
should overestimate or underestimate.

\begin{center}
\renewcommand{\arraystretch}{1.7}
\begin{tabular}{p{0.30\textwidth}|p{0.54\textwidth}}
\textbf{shape condition} & \textbf{consequence} \\ \hline

\(f\) increasing on \([a,b]\) &
\(\displaystyle L_n\le \int_a^b f(x)\,dx\le R_n\) \\

\(f\) decreasing on \([a,b]\) &
\(\displaystyle R_n\le \int_a^b f(x)\,dx\le L_n\) \\

\(f\) concave up on \([a,b]\) &
The trapezoidal rule overestimates. \\

\(f\) concave down on \([a,b]\) &
The trapezoidal rule underestimates.
\end{tabular}
\end{center}

For the midpoint rule, the concavity effect is opposite:

\begin{center}
\renewcommand{\arraystretch}{1.6}
\begin{tabular}{c|c}
\textbf{condition} & \textbf{midpoint rule behavior} \\ \hline
concave up & midpoint rule underestimates \\
concave down & midpoint rule overestimates
\end{tabular}
\end{center}

\subsection{Choosing a numerical rule}
\label{appsubsec:choosing-numerical-rule}

\begin{center}
\renewcommand{\arraystretch}{1.7}
\begin{tabular}{p{0.26\textwidth}|p{0.58\textwidth}}
\textbf{situation} & \textbf{reasonable choice} \\ \hline

Only left-endpoint data are available &
Left rectangle rule \\

Only right-endpoint data are available &
Right rectangle rule \\

Midpoint data are available or easy to compute &
Midpoint rule \\

Endpoint data are available and graph is nearly linear on small intervals &
Trapezoidal rule \\

Function values are easy to compute and high accuracy is desired &
Simpson's rule \\

A guaranteed error tolerance is required &
Use a rule with an error bound and choose \(n\) from the bound.
\end{tabular}
\end{center}

\begin{quote}
\textbf{Final numerical habit.}
A decimal approximation should come with a reason to trust it:

\[
\text{shape check, error bound, or convergence evidence.}
\]

Best of all, use more than one.
\end{quote}